\chapter{Additional CVEs Investigated}

During the research, three additional eBPF verifier CVEs were analyzed but not included in the main body of this report due to practical limitations encountered during exploitation. Each investigation contributed to a deeper understanding of the verifier's defense mechanisms and informed the strategies used in the primary CVEs.

\section{CVE-2021-33200 — \texttt{alu\_limit} Masking Direction Bypass}

\begin{table}[H]
\centering
\begin{tabular}{@{}ll@{}}
\toprule
\textbf{Component} & \textbf{Details} \\
\midrule
Bug class & Logic error in pointer arithmetic sanitization \\
Component & \texttt{sanitize\_ptr\_alu()} in \texttt{verifier.c} \\
Kernel & 5.10.15 \\
Severity & 7.8 HIGH \\
\bottomrule
\end{tabular}
\end{table}

\textbf{Root Cause:} When an offset register's value range crosses zero (\texttt{smin\_value < 0 < smax\_value}), the \texttt{alu\_limit} masking direction computation produces an overly permissive limit, allowing OOB pointer arithmetic even for unprivileged programs. This CVE directly targets the \texttt{alu\_limit} defense that blocked our CVE-2021-3600 LPE attempt.

\textbf{Status:} Milestones 1--4 were implemented and tested. The verifier accepted programs with zero-crossing offsets as expected, but achieving reliable OOB access required precise control of register bounds across the zero-crossing boundary. The exploitation chain was not completed due to the complexity of the masking direction change and the time required to craft stable bounds.

\textbf{Academic Value:} Understanding this CVE was essential for the project's narrative—it demonstrates that \texttt{alu\_limit} itself had vulnerabilities, providing a bridge between CVE-2021-3600 (blocked by \texttt{alu\_limit}) and CVE-2021-3444 (exploited on a kernel without \texttt{alu\_limit}).


\section{CVE-2023-52462 — Spilled Pointer Corruption}

\begin{table}[H]
\centering
\begin{tabular}{@{}ll@{}}
\toprule
\textbf{Component} & \textbf{Details} \\
\midrule
Bug class & Spilled pointer corruption check bypass \\
Component & Stack-slot tracking in \texttt{verifier.c} \\
Kernel & 6.7.1 \\
Severity & 5.5 MEDIUM \\
\bottomrule
\end{tabular}
\end{table}

\textbf{Root Cause:} The verifier's spilled-pointer corruption checks use an index that can be bypassed after a crafted partial byte write. This allows partial pointer corruption that should be rejected, enabling pointer-value disclosure.

\textbf{Status:} Three milestones were completed: the corruption sequence was triggered (M1), kernel pointer values were leaked for KASLR bypass (M2), and reliability was assessed (M3). However, direct type-confusion dereference of reloaded corrupted values was blocked because the verifier correctly downgrades reloaded state to scalar.

\textbf{Why Stopped:} The primary impact is limited to information disclosure. Achieving arbitrary R/W from this primitive requires chaining with an independent memory-corruption vulnerability, which was beyond the scope of this single-CVE analysis.


\section{CVE-2023-52676 — 32-bit Integer Overflow in Stack Bounds}

\begin{table}[H]
\centering
\begin{tabular}{@{}ll@{}}
\toprule
\textbf{Component} & \textbf{Details} \\
\midrule
Bug class & 32-bit integer overflow in stack bounds arithmetic \\
Component & \texttt{check\_stack\_access\_within\_bounds()} in \texttt{verifier.c} \\
Kernel & 6.7.1 \\
Severity & 5.5 MEDIUM \\
\bottomrule
\end{tabular}
\end{table}

\textbf{Root Cause:} The function \texttt{check\_stack\_access\_within\_bounds()} computes \texttt{int min\_off = reg->var\_off.value + off}, where both operands can be large. The sum overflows \texttt{int}, wrapping to a small value and causing the verifier to incorrectly approve OOB stack accesses.

\textbf{Status:} A comprehensive defense-in-depth analysis was performed across 5 exploitation strategies. All strategies were blocked by layered defenses: \texttt{adjust\_ptr\_min\_max\_vals()} rejects large constants before they reach the vulnerable code, and the overflow is only triggerable through kfunc argument validation paths requiring \texttt{CAP\_BPF} and BTF support.

\textbf{Why Stopped:} The vulnerability is real but the attack surface is restricted to privileged contexts with kfunc access (\texttt{BPF\_PROG\_TYPE\_TRACING}). No exploitation path was found for unprivileged programs, and the defense-in-depth analysis was documented as the primary research output.


\section{Key Takeaways from Additional Investigations}

\begin{enumerate}
    \item \textbf{Defense-in-depth works:} Both CVE-2023-52676 and CVE-2023-52462 demonstrated that even when a specific verifier check is flawed, additional layers (pointer type restrictions, stack bounds guards, scalar downgrades) can independently prevent exploitation.
    \item \textbf{Privilege boundaries matter:} Several CVEs are only reachable with specific capabilities (\texttt{CAP\_BPF}, \texttt{CAP\_NET\_ADMIN}), which significantly reduces their real-world impact for unprivileged LPE.
    \item \textbf{Failed approaches inform success:} The analysis of CVE-2021-33200's \texttt{alu\_limit} bypass directly motivated the decision to target CVE-2021-3444 on kernel 4.19.19, where \texttt{alu\_limit} is absent entirely—leading to the project's successful full LPE.
\end{enumerate}
